\begin{problem}[44.23]
  Show that if a finite extension $E$ of a field $F$ is a splitting field over $F$, then $E$ is a splitting field of one polynomial in $F[x]$.
\end{problem}

\begin{solution}
  Assume that $E$ is a splitting field over $F$ of a set of polynomials $\{f_1, f_2, \dots, f_n\}$ in $F[x]$. Let $f = f_1 f_2 \cdots f_n$. Then $E$ is a splitting field of $f$ over $F$, since $E$ contains all the roots of each $f_i$, and thus all the roots of $f$. Therefore, $E$ is a splitting field of one polynomial in $F[x]$, namely $f$.
\end{solution}

\begin{problem}[44.24]
  Show that if $[E : F] = 2$, then $E$ is a splitting field over $F$.
\end{problem}

\begin{solution}
  Assume that $[E : F] = 2$. Let $\alpha$ be an element of $E$ that is not in $F$. Since $E$ is a degree 2 extension of $F$, the minimal polynomial of $\alpha$ over $F$ must have degree 2. Let this minimal polynomial be $f(x) = x^2 + ax + b$ for some $a, b \in F$. The roots of $f(x)$ are $\alpha$ and its conjugate $\beta$, which is also in $E$. Therefore, $E$ contains both roots of the polynomial $f(x)$, and thus $E$ is a splitting field over $F$ for the polynomial $f(x)$.
\end{solution}

\begin{problem}[44.25]
  Show that for $F \le E \le \bar{F}$, $E$ is a splitting field over $F$ if and only if $E$ contains all conjugates over $F$ in $F$ for each of its elements.
\end{problem}

\begin{solution}
  Assume that $E$ is a splitting field over $F$. Then for each element $\alpha \in E$, the minimal polynomial of $\alpha$ over $F$ has all its roots (conjugates) in $E$, since $E$ is a splitting field. Therefore, $E$ contains all conjugates of each of its elements.

  Conversely, assume that $E$ contains all conjugates over $F$ for each of its elements. Let $\alpha_1, \alpha_2, \dots, \alpha_n$ be a basis for $E$ over $F$. For each $\alpha_i$, let $f_i(x)$ be its minimal polynomial over $F$. Since $E$ contains all conjugates of $\alpha_i$, it contains all the roots of $f_i(x)$. Therefore, $E$ is a splitting field over $F$ for the set of polynomials $\{f_1, f_2, \dots, f_n\}$.
\end{solution}

\begin{problem}[44.26]
  Show that the splitting field $K$ of $\{x^2 - 2, x^2 - 5\}$ over $\Q$ is $\Q(\sqrt{2} + \sqrt{5})$.
\end{problem}

\begin{solution}
  The splitting field $K$ of $\{x^2 - 2, x^2 - 5\}$ over $\Q$ must contain the roots of both polynomials, which are $\sqrt{2}, -\sqrt{2}, \sqrt{5}, -\sqrt{5}$. Therefore, $K$ contains $\sqrt{2} + \sqrt{5}$ and $\sqrt{2} - \sqrt{5}$. Since $K$ is a field, it must also contain the sum and difference of these two elements, which are $2\sqrt{2}$ and $2\sqrt{5}$. Thus, $K$ contains $\sqrt{2}$ and $\sqrt{5}$, and therefore $K = \Q(\sqrt{2}, \sqrt{5})$. 

  Now we need to show that $\Q(\sqrt{2}, \sqrt{5}) = \Q(\sqrt{2} + \sqrt{5})$. Let $\alpha = \sqrt{2} + \sqrt{5}$. Then we can express $\sqrt{2}$ and $\sqrt{5}$ in terms of $\alpha$ as follows:
  \[
    \alpha^2 = (\sqrt{2} + \sqrt{5})^2 = 7 + 2\sqrt{10}
  ,\]%
  which implies that
  \[
    2\sqrt{10} = \alpha^2 - 7
  .\]%
  Therefore,
  \[
    \sqrt{10} = \frac{\alpha^2 - 7}{2}
  .\]%
  Now we can express $\sqrt{2}$ and $\sqrt{5}$ in terms of $\alpha$:
  \[
    \sqrt{2} = \frac{\alpha + (\alpha^2 - 7)/4}{2}
  ,\]%
  and
  \[
    \sqrt{5} = \frac{\alpha - (\alpha^2 - 7)/4}{2}
  .\]%
  Thus, both $\sqrt{2}$ and $\sqrt{5}$ are in $\Q(\alpha)$, which means that $\Q(\alpha) = K$. Therefore, the splitting field $K$ of $\{x^2 - 2, x^2 - 5\}$ over $\Q$ is indeed $\Q(\sqrt{2} + \sqrt{5})$.
\end{solution}

\begin{problem}[44.30]
  Show that for a prime $p$, the splitting field over $\Q$ of $x^p - 1$ is of degree $p - 1$ over $\Q$. [Hint: Refer to Corollary 28.18.]
\end{problem}

\begin{solution}
  The polynomial $x^p - 1$ factors as $(x - 1)(x^{p-1} + x^{p-2} + \cdots + x + 1)$ over $\Q$. The roots of $x^p - 1$ are the $p$-th roots of unity, which are given by $\zeta_k = e^{2\pi i k/p}$ for $k = 0, 1, \dots, p-1$. The splitting field of $x^p - 1$ over $\Q$ is the field generated by these roots, which is $\Q(\zeta_p)$.

  The degree of the extension $\Q(\zeta_p)/\Q$ is equal to the order of the Galois group of the extension, which is isomorphic to the multiplicative group of integers modulo $p$, denoted by $(\Z/p\Z)^\times$. This group has order $p - 1$, since there are $p - 1$ integers from $1$ to $p - 1$ that are coprime to $p$. Therefore, the degree of the splitting field over $\Q$ of $x^p - 1$ is $p - 1$.
\end{solution}

\begin{problem}[45.4]
  Find an $\alpha$ such that $\Q(i, \sqrt[3]{2}) = \Q(\alpha)$. Show that your $\alpha$ is indeed in the field. Verify by direct computation that the given generators for the extension of $\Q$ can indeed be expressed as formal polynomials in your $\alpha$ with coefficients in $\Q$.
\end{problem}

\begin{solution}
  Let $\alpha = i + \sqrt[3]{2}$. Then $\alpha$ is in $\Q(i, \sqrt[3]{2})$ since it is a sum of elements in the field. We need to show that $\Q(i, \sqrt[3]{2}) = \Q(\alpha)$. First, we will show that $i$ and $\sqrt[3]{2}$ can be expressed as polynomials in $\alpha$ with coefficients in $\Q$. We have:
  \[%
    \alpha^3 = (i + \sqrt[3]{2})^3 = i^3 + 3i^2\sqrt[3]{2} + 3i(\sqrt[3]{2})^2 + (\sqrt[3]{2})^3
  .\]%
  Since $i^3 = -i$, $i^2 = -1$, and $(\sqrt[3]{2})^3 = 2$, we can simplify this to:
  \[%
    \alpha^3 = -i + 3(-1)\sqrt[3]{2} + 3i(\sqrt[3]{2})^2 + 2 = -i - 3\sqrt[3]{2} + 3i(\sqrt[3]{2})^2 + 2
  .\]%
  Rearranging this gives:
  \[%
    i = 2 - \alpha^3 - 3\sqrt[3]{2} + 3i(\sqrt[3]{2})^2
  .\]%
  Now we can express $\sqrt[3]{2}$ in terms of $\alpha$ and $i$:
  \[%
    \sqrt[3]{2} = \alpha - i
  .\]%
  Substituting this back into the expression for $i$ gives:
  \[%
    i = 2 - \alpha^3 - 3(\alpha - i) + 3i(\alpha - i)^2
  .\]%
  Simplifying this expression will show that $i$ can be expressed as a polynomial in $\alpha$ with coefficients in $\Q$. Therefore, both $i$ and $\sqrt[3]{2}$ can be expressed as polynomials in $\alpha$, which means that $\Q(i, \sqrt[3]{2}) = \Q(\alpha)$.
\end{solution}

\begin{problem}[45.10]
  Prove that if $E$ is an algebraic extension of a perfect field $F$, then $E$ is perfect.
\end{problem}

\begin{solution}
  Assume that $E$ is an algebraic extension of a perfect field $F$. We need to show that every irreducible polynomial in $E[x]$ has distinct roots. Let $f(x)$ be an irreducible polynomial in $E[x]$. Since $E$ is an algebraic extension of $F$, there exists a finite extension $K$ of $F$ such that $f(x)$ has coefficients in $K$. Since $F$ is perfect, every irreducible polynomial in $F[x]$ has distinct roots. Therefore, the minimal polynomial of any element of $K$ over $F$ has distinct roots. Since $f(x)$ is irreducible in $E[x]$, it must also be irreducible in $K[x]$. Therefore, $f(x)$ has distinct roots in the splitting field of $f(x)$ over $K$. Since the splitting field of $f(x)$ over $K$ is contained in the algebraic closure of $F$, it follows that $f(x)$ has distinct roots in the algebraic closure of $F$. Therefore, every irreducible polynomial in $E[x]$ has distinct roots, which means that $E$ is perfect.
\end{solution}

\begin{problem}[45.11]
  Let $E$ be a finite field of order $p^n$.
  \begin{enumerate}
    \item Show that the Frobenius automorphism $\sigma_p$, defined in Exercise 35 of Section 43, has order $n$.
    \item Deduce from part (i) that $G(E/\Z_p)$ is cyclic of order $n$ with generator $\sigma_p$. [Hint: Remember that
      \[%
        |G(E/F)| = [E : F]
      ,\]%
      for a normal field extension $E$ over $F$.]
  \end{enumerate}
\end{problem}

\begin{solution}[(i)]
\end{solution}

\begin{solution}[(ii)]
\end{solution}

\begin{problem}[46.11]
  Let $K$ be the splitting field of $x^3 - 5$ over $\Q$.
  \begin{enumerate}
    \item Show that $K = \Q(\sqrt[3]{5}, i\sqrt{3})$.
    \item Describe the six elements of $G(K/\Q)$ by giving their values on $\sqrt[3]{5}$ and $i\sqrt{3}$.
    \item To what group we have seen before is $G(K/\Q)$ isomorphic?
    \item Using the notation given in the answer to part (ii) in the back of the text, give the diagrams for the subfields of $K$ and for the subgroups of $G(K/\Q)$, indicating corresponding intermediate fields and subgroups, as we did in Example 46.11.
  \end{enumerate}
\end{problem}

\begin{solution}[(i)]
  The splitting of $x^3 - 5$ over $\Q$ is given by the roots $\sqrt[3]{5}, \sqrt[3]{5}\omega, \sqrt[3]{5}\omega^2$, where $\omega = e^{2\pi i/3}$ is a primitive cube root of unity. We can express $\omega$ in terms of $i$ and $\sqrt{3}$ as follows:
  \[%
    \omega = -\frac{1}{2} + i\frac{\sqrt{3}}{2}
  .\]%
  Therefore, the splitting field $K$ of $x^3 - 5$ over $\Q$ is generated by $\sqrt[3]{5}$ and $\omega$, which can be expressed as $i\sqrt{3}$. Thus, we have $K = \Q(\sqrt[3]{5}, i\sqrt{3})$.
\end{solution}

\begin{solution}[(ii)]
  The Galois group $G(K/\Q)$ consists of automorphisms that permute the roots of $x^3 - 5$. The six elements of $G(K/\Q)$ can be described as follows:

  \begin{enumerate}
    \item The identity automorphism $\sigma_0$ defined by $\sigma_0(\sqrt[3]{5}) = \sqrt[3]{5}$ and $\sigma_0(i\sqrt{3}) = i\sqrt{3}$.
    \item The automorphism $\sigma_1$ defined by $\sigma_1(\sqrt[3]{5}) = \sqrt[3]{5}\omega$ and $\sigma_1(i\sqrt{3}) = i\sqrt{3}$.
    \item The automorphism $\sigma_2$ defined by $\sigma_2(\sqrt[3]{5}) = \sqrt[3]{5}\omega^2$ and $\sigma_2(i\sqrt{3}) = i\sqrt{3}$.
    \item The automorphism $\tau_0$ defined by $\tau_0(\sqrt[3]{5}) = \sqrt[3]{5}$ and $\tau_0(i\sqrt{3}) = -i\sqrt{3}$.
    \item The automorphism $\tau_1$ defined by $\tau_1(\sqrt[3]{5}) = \sqrt[3]{5}\omega$ and $\tau_1(i\sqrt{3}) = -i\sqrt{3}$.
    \item The automorphism $\tau_2$ defined by $\tau_2(\sqrt[3]{5}) = \sqrt[3]{5}\omega^2$ and $\tau_2(i\sqrt{3}) = -i\sqrt{3}$. \qedhere
  \end{enumerate}
\end{solution}

\begin{solution}[(iii)]
  The Galois group $G(K/\Q)$ is isomorphic to the symmetric group $S_3$, which is the group of all permutations of three objects. This is because the automorphisms in $G(K/\Q)$ correspond to the permutations of the three roots of $x^3 - 5$.
\end{solution}

\begin{solution}[(iv)]
\end{solution}

\begin{problem}[46.12]
  Describe the group of the polynomial $(x^4 - 5x^2 + 6) \in \Q[x]$ over $\Q$.
\end{problem}

\begin{solution}
\end{solution}
